\subsection{World model, planner, and optional policy}
\label{sec:planner-actions}
An action-conditioned world model predicts consequences; a planner selects
actions; a policy network amortizes that selection. These are distinct objects.
For partially observed deployments, let $h_t$ contain only observations available
by decision time, their ages, and previous applied actions. The proposed
transition model is $F_\theta(\cdot\mid h_t,\bar a_t,c_T,\Delta_t)$, where
$\bar a_t$ is the action that actually takes effect and $\Delta_t$ is its
duration. A requested action can be rejected, delayed, or clipped, and must not
silently replace $\bar a_t$ in a training record. A receding-horizon planner would
solve
\begin{equation}
 A^* = \arg\max_{A\in\mathcal A(h_t)}
 \E_{F_\theta}\!\left[\sum_{k=0}^{H_p-1}\gamma^k
 U(\hat s_{t+k+1},a_{t+k};w_T)\right],
 \label{eq:planner-contract}
\end{equation}
Here $H_p$ is the planning horizon, $\gamma\in[0,1]$ a discount factor, and
$\mathcal A(h_t)$ the feasible candidate-sequence set. The search respects
actuator limits, service constraints, and a fixed computation budget.
Only its first accepted action is executed before replanning from new evidence.
Finite enumeration, sampling, or model predictive control can implement this
planner without training an actor. An optional actor
$\pi_\phi(a\mid h_t,w_T)$ can later be trained from evaluated planner decisions
or by reinforcement learning. Historical operator actions are behavior samples,
not automatically optimal-action labels.

Equation~\eqref{eq:planner-contract} is a deployment contract, not a reported
closed-loop result. The implemented \texttt{planning\_audit} ranks a finite grid
of two-link power vectors by \emph{one-step} utility. It repeats each candidate
over the forecast call but scores only its first predicted state; it neither
searches general action sequences nor trains an actor.

\subsection{A bounded, hierarchical RAN action space}
Table~\ref{tab:planner-actions} identifies actuators rather than listing KPIs
as actions. A slow controller can choose slice budgets, scheduler modes, and
power envelopes; a local controller realizes the corresponding fast allocations.
The same actuator must have one owner per control interval. The planner may
change only fields exposed by the selected simulator or execution interface;
capability masks disable unsupported actions. Neither the slow language backbone
nor a nominal model step implies permission to perform slot-level scheduling
through a Near-RT RIC. Placement and deadlines require a separate integration
measurement~\cite{oran2026}.

\begin{table*}[t]
\centering
\small
\caption{Proposed RAN action contract. Only the first row is exercised in the
current synthetic decision audit. The other rows define scoped extensions,
not implemented or experimentally validated capabilities. Rates describe
ownership and dependencies, not measured execution deadlines.}
\label{tab:planner-actions}
\begin{tabular}{p{0.15\textwidth}p{0.27\textwidth}p{0.25\textwidth}p{0.23\textwidth}}
\toprule
Action family & Decision variables & Rate and required evidence & Essential limits \\
\midrule
Power / interference & Per-link power or power envelope; current reference uses
two normalized powers & Local fast adjustment under slower envelopes;
applied power, channel/interference, load & Total/per-link power, ramp rate,
neighbor interference \\
Resources / scheduling & Slice PRB or RBG budget; scheduler mode; optionally local
UE allocations & Budgets/modes at a supervisory rate; local scheduling faster;
queue, traffic, CQI, granted resources & Resource conservation, minimum service,
integer allocations, one actuator owner \\
Beam management & Codebook beam, candidate set, sweep/measurement trigger &
Channel/blockage dependent; beam powers, alignment, switching feedback & Valid
codebook, switching and measurement overhead \\
Mobility / association & Target cell, handover trigger or parameters, association &
Event driven; neighbor measurements and handover outcomes & Dwell/hysteresis,
interruption cost, reachable target \\
Admission / energy & Admit or defer traffic; carrier/cell sleep and wake &
Load driven, with slower activation effects; demand and energy measurements &
Coverage, wake delay, service obligations \\
Measurement / reporting & Measurement period, report rate, sensor activation &
Information refresh at multiple rates; observation ages, acquisition cost &
Maximum staleness and acquisition budget; model the changed observation process \\
\bottomrule
\end{tabular}
\end{table*}

Throughput, delay, and energy are outcomes or objective terms, not directly
commandable actions. External traffic, mobility, and channel evolution are
exogenous unless a particular task explicitly controls them. Objective weights
are policy configuration, not physical interventions. Selecting a model,
refreshing context, or requesting replanning is a supervisory operation; it is
not evidence that a radio action was applied. Edge offloading and cache placement
would require additional compute/cache state and actuators and are outside this
paper's evaluated action space. Every extension also needs an explicit
hold/fallback behavior; zero transmit power is not a universal safe fallback.

\subsection{What public data can supervise}
\label{sec:action-data}
The absence of action logs in UC6 and UCC production traces does not imply that
all public wireless datasets lack action information. Three evidence types
must be separated.

\paragraph{Passive commercial observations.}
UC6 and UCC support observation forecasting, as evaluated in
Section~\ref{sec:uc6}. They do not provide the applied RAN control and effective
time needed to identify power or resource interventions. Unknown actions must
be masked rather than encoded as zero power.

\paragraph{Measured alternatives for beam selection.}
DeepSense task releases include measured beam-power vectors and best-beam
indices derived from them~\cite{deepsensebeam}. These can supervise beam
selection and, with a matched evaluation target, ranking loss relative to the
measured best beam. An index derived by taking an argmax is not an executed
control trajectory. Beam sweeps alone do not identify queue evolution,
handover interruption, or the long-term effects of changing beams. The measured
power vector must be a target rather than an unavailable decision-time input.

\paragraph{Configuration-conditioned RF-emulator experiments.}
The public Colosseum COMMAG and ColO-RAN releases document slicing allocations
and RR/WF/PF scheduling configurations~\cite{commagdata,colorandata}.
Our source audit on 29 October 2026 downloaded one per-UE CSV from each release
and found timestamp, \texttt{slice\_prb}, \texttt{scheduling\_policy},
\texttt{power\_multiplier}, queue, CQI, throughput, and requested/granted-resource
fields. Both inspected files contain a single resource configuration,
scheduler mode, and power multiplier; a column's presence is not evidence of
intervention coverage. The COMMAG README describes a resizing schedule, but
that schedule cannot substitute for changes observed in a particular CSV.
The source revisions, sample paths, hashes, and findings are recorded in the
repository's action-data audit.

Commercial Traffic Twinning provides a complementary collection: a
one-BS/eight-UE Colosseum experiment replays commercial traffic under five
slice-resource configurations and two scheduling modes, with timestamped
PHY/MAC/application measurements~\cite{twinning2024}. Those interventions and
feedback occur in the RF emulator; the commercial origin of the offered traffic
does not make the resulting action-response trajectories production-network
interventions. Its published MAC counters aggregate over 250~ms and do not
recover individual slot allocations. This audit inspected its documentation,
not its full data archives.

These releases are candidates for \emph{restricted}, configuration-conditioned
training after entity, unit, timing, coverage, and license checks. Constant-action
runs can teach within-configuration responses, but do not alone identify
switching transients or arbitrary multistep control. Resource demand and grants
aggregated during an interval are not automatically decision-time actions.
Trial identifiers or future KPI windows must not leak into planner inputs.
None of these releases has been used to train the models reported here, and
none establishes a complete, jointly observed action space for Table~\ref{tab:planner-actions}.

\subsection{Closing action-coverage gaps with controlled simulation}
For the broad sequential-control extension, our proposed data route combines
measured representations and restricted public control experiments with
\emph{explicitly labeled} simulator interventions. Simulation is not necessary
for every task: measured beam-selection supervision already exists. It is a
practical source of missing action coverage and switching outcomes when
corresponding controlled measurements are unavailable.

The first extension should retain power control as a synthetic regression
test and separately study slice budgets plus scheduler selection in a
packet-level environment. It should not claim simultaneous validation of
beams, handovers, energy, and admission. ns-3/5G-LENA offers channel, queue,
scheduler, traffic, and power-control components~\cite{lena2026}; the chosen
version still needs task-specific actuator hooks and logging. A propagation
or ray-tracing model alone does not generate scheduler/queue feedback.

Each episode should preserve available observation history, exogenous inputs,
requested and applied actions, effective times and durations, execution status,
subsequent independently computed KPI/cost measurements, behavior policy,
simulator version, seeds, and provenance. Collect rule-based and varied
feasible policies, not only an expert's preferred actions. From a saved simulator
state, replay alternatives with matched exogenous realizations to measure action
effects, retaining the shared branch origin for grouped train/test splitting.
Fit action-conditioned dynamics first; evaluate a bounded search planner next;
distill an actor only if search latency justifies it. Report multistep error,
action-ranking regret, realized return, service violations, switching costs,
and total decision latency against rules and an action-independent predictor.
Use held-out scenarios and independently acquired testbed/field feedback before
claiming transfer. LLMs may propose traffic scenarios or constraints, but the
simulator or instrumented network must compute physical outcomes. Synthetic,
RF-emulator, derived-label, and field evidence remain separate throughout.
